\documentclass{article}

\usepackage[italian]{babel}
\usepackage[margin=20mm, footskip=20pt]{geometry}
\usepackage{graphicx}
\usepackage{subfiles}
\usepackage{hyperref}
\usepackage{nameref}
\usepackage{titlesec}
\usepackage{longtable}
\usepackage[table]{xcolor}
\usepackage{titling}
\usepackage{lastpage}
\usepackage{ifthen}
\usepackage{calc}
\usepackage{soulutf8}
\usepackage{contour}
\usepackage{float}
\usepackage{fancyhdr}
\usepackage{multirow}
\usepackage{pgfgantt}
\usepackage{lscape}
\usepackage{background}
\usepackage{lmodern}
\usepackage{textcomp}
\usepackage{lastpage}
\usepackage[utf8]{inputenc}
\usepackage{makecell}

% definizione dei percorsi in cui cercare immagini
\graphicspath{ {./}
    {./src/img/}
}

% setup della sottolineatura
\setuldepth{Flat}
\contourlength{0.8pt}

\newcommand{\uline}[1]{%
  \ul{{\phantom{#1}}}%
  \llap{\contour{white}{#1}}%
}


% setup dei link
\hypersetup{
  % set true if you want colored links (instead of boxes)
  colorlinks=true,
  % set to all if you want both sections and subsections linked
  linktoc=all,
  % set color for file links
  filecolor=blue,
  % set color for internal links
  linkcolor=black,
  % set url color
  urlcolor=blue,
  % set characters encoding in the bookmarks tab
  pdfencoding=unicode,
}

% setup forma \paragraph e \subparagraph
\titleformat{\paragraph}[hang]{\normalfont\normalsize\bfseries}{\theparagraph}{1em}{}
\titleformat{\subparagraph}[hang]{\normalfont\normalsize\bfseries}{\thesubparagraph}{1em}{}

% setup profondità indice di default
\setcounter{secnumdepth}{5}
\setcounter{tocdepth}{5}

\makeatletter %% non togliere, i comandi che definiscono i placeholder vanno qui
% esempio di utilizzo: \appendToGraphicspath{./img/} (un comando diverso per ogni path da includere)
% N.B.: ci DEVE essere un forward slash alla fine del path, a indicare che è una cartella.
\newcommand\appendToGraphicspath[1]{%
  \g@addto@macro\Ginput@path{{#1}}%
}

\newcommand{\setTitle}[1]{%
  \newcommand{\@placeholderTitle}{#1}%
}
\newcommand{\placeholderTitle}{\@placeholderTitle}

\newcommand{\setUso}[1]{%
  \newcommand{\@uso}{#1}%
}
\newcommand{\uso}{\@uso}

\newcommand{\setVersione}[1]{%
  \newcommand{\@versione}{#1}%
}
\newcommand{\versione}{\@versione}

\newcommand{\disabilitaVersione}{%
  \renewcommand{\setVersione}[1]{}%
  \renewcommand{\versione}{DISABILITATA}
}

\newcommand{\setResponsabile}[1]{%
  \newcommand{\@responsabile}{#1}%
}
\newcommand{\responsabile}{\@responsabile}

\newcommand{\setRedattori}[1]{%
  \newcommand{\@redattori}{#1}%
}
\newcommand{\redattori}{\@redattori}

\newcommand{\setVerificatori}[1]{%
  \newcommand{\@verificatori}{#1}%
}
\newcommand{\verificatori}{\@verificatori}

\newcommand{\setDestinatari}[1]{%
  \newcommand{\@destinatari}{#1}%
}
\newcommand{\destinatari}{\@destinatari}

\newcommand{\setDescrizione}[1]{%
  \newcommand{\@descrizione}{#1}%
}
\newcommand{\descrizione}{\@descrizione}

\newcommand{\setModifiche}[1]{%
  \newcommand{\@modifiche}{#1}%
}

\newcommand{\modifiche}{\@modifiche}
\makeatother %% non togliere, i comandi che definiscono i placeholder vanno qui

% hook per lo script che genera il glossario
\newcommand{\glossario}[1]{#1\textsubscript{G}}

% comandi per rendere opzionali gli elenchi di figure
\newcommand{\elencoFigure}{%
  \renewcommand{\listfigurename}{Elenco delle figure}%
  \listoffigures%
}

\newcommand{\disabilitaElencoFigure}{%
  \renewcommand{\elencoFigure}{}%
}

% comandi per rendere opzionali le tabelle
\newcommand{\elencoTabelle}{%
  \renewcommand{\listtablename}{Elenco delle tabelle}%
  \listoftables%
}

\newcommand{\disabilitaElencoTabelle}{%
  \renewcommand{\elencoTabelle}{}%
}

% comando per riferirsi ad una sezione
\newcommand{\refSec}[1]{%
  \hyperref[#1]{\S\ref*{#1}}%
}


\appendToGraphicspath{../../config/src/img/}

\setTitle{Verbale Esterno 2021-02-17}

\setVersione{v0.2.0}

\setResponsabile{
	Pagotto Manuel
}

\setRedattori{
	Contro Daniel Eduardo
}

\setVerificatori{
	Paparazzo Giorgia
}

\setUso{Esterno}

\setDestinatari{
	prof. Vardanega Tullio \\ &
	prof. Cardin Riccardo \\ &
	SpaghettiCode
}

\setDescrizione{Riassunto dell'incontro tenutosi il 17 febbraio 2021 realizzato dal gruppo SpaghettiCode in forma di meeting online con Piccoli Gregorio.}

\setModifiche{
	v0.2.0 & Contro Daniel Eduardo & Amministratore & 2021-02-27 & Corretto \refSec{sec:resoconto} \\
	v0.1.0 & Contro Daniel Eduardo & Amministratore & 2021-02-26 & Modifiche a \refSec{sec:resoconto} \\
	v0.0.1 & Contro Daniel Eduardo & Amministratore & 2021-02-25 & Creazione del documento e prima stesura
}

\disabilitaElencoFigure
\disabilitaElencoTabelle

\begin{document}

\pagenumbering{gobble}


% variabile prima pagina (serve nell'istruzione dopo per stampare il bg solo nella prima pagina)
\newif\iffirstpage
\firstpagetrue

% setta l'immagina di bg
\backgroundsetup{
	scale=1,
	opacity=0.2,
	placement=top,
	contents={%
			\iffirstpage
				\includegraphics[width=\paperwidth]{datascience_og_colori.png}%
				\global\firstpagefalse
			\fi
		}%
}

\begin{titlepage}% per non stampare il numero della pagina

	\centering % allinea al centro la pagina
	\hspace{0.05\textwidth}% spazio tra linea e testo
	% lasciare questa riga per il corretto funzionamento di \parbox
	\parbox[b]{0.4\textwidth}{% cambiando la larghezza del testo il paragrafo si muove a destra o a sinistra 
	{\hspace{0.05\textwidth}\includegraphics[width=4cm,height=4cm]{logo_colori.png}}\\[3\baselineskip] % logo
	{\Huge\bfseries SpaghettiCode}\\ [\baselineskip] %titolo
	{\texttt{spaghetti.code.g6@gmail.com}}\\[\baselineskip]\\[4\baselineskip] % 
	{\Large\textsc\mbox{\placeholderTitle{}}}\\[4\baselineskip] % nome del documento
	{\begin{tabular}{r|l}
		\hline                                  \\
		% testo in grassetto
		\textbf{Versione}     & \versione{}     \\
		\rule{0pt}{3ex}%  EXTRA vertical height 
		\textbf{Approvazione} & \responsabile{} \\
		\rule{0pt}{3ex}%  EXTRA vertical height 
		\textbf{Redazione}    & \redattori{}    \\
		\rule{0pt}{3ex}%  EXTRA vertical height 
		\textbf{Verifica}     & \verificatori{} \\
		\rule{0pt}{3ex}%  EXTRA vertical height 
		\textbf{Uso}          & \uso{}          \\
		\rule{0pt}{3ex}%  EXTRA vertical height 
		\textbf{Destinato a}  & \destinatari{}
		\ifthenelse{\equal{\uso}{Esterno}}{
		\\ & Zucchetti S.p.A.
		}{}
	\end{tabular}}\\[4\baselineskip]

	}

	{\bfseries Descrizione}\\
	{\descrizione{}}\\[1\baselineskip]



\end{titlepage}

\newgeometry{textheight=660pt, lmargin=2cm, tmargin=2cm, rmargin=2cm}

% setup di header e footer nelle pagine senza numero
\fancypagestyle{nopage}{%
	\fancyhf{}%
	\fancyhead[R]{\includegraphics[width=1.3cm]{logo_colori.png}}%
	\fancyhead[L]{\emph{SpaghettiCode}\\\placeholderTitle{}}%
}
% setup di header e footer nelle pagine col numero
\fancypagestyle{usual}{%
	\fancyhf{}%
	\fancyhead[R]{\includegraphics[width=1.3cm]{logo_colori.png}}%
	\fancyhead[L]{\emph{SpaghettiCode}\\\placeholderTitle{}}%
	\fancyfoot[R]{\thepage\ di~\pageref{LastPage}}%
}
\setlength{\headheight}{1.8cm}

\newpage
\pagestyle{nopage}

\setcounter{table}{-1}

%REGISTRO DELLE MODIFICHE

\section*{Registro delle modifiche}
\label{sec:registro_delle_modifiche}

\rowcolors{2}{white!80!lightgray!90}{white}
\renewcommand{\arraystretch}{2} % allarga le righe con dello spazio sotto e sopra

\begin{longtable}
	[H]{|>{\centering\bfseries}m{2cm}|>{\centering}m{3.5cm}|>{\centering}m{2.5cm}|>{\centering}m{3cm}|>{\centering\arraybackslash}m{5cm}|}
	
	\hline
	\rowcolor{lightgray}
	{\textbf{Versione}} & {\textbf{Nominativo}} & {\textbf{Ruolo}} & {\textbf{Data}} & {\textbf{Descrizione}} \\
	\hline
	\endfirsthead
	
	\hline
	\rowcolor{lightgray}
	{\textbf{Versione}} & {\textbf{Nominativo}} & {\textbf{Ruolo}} & {\textbf{Data}} & {\textbf{Descrizione}} \\
	\hline
	\endhead
	
	\hline
	\rowcolor{white}
 	\multicolumn{5}{|c|}{\emph{Continua alla pagina successiva...}}\\
	\hline

	\endfoot
	\hline
	\endlastfoot
	
	\modifiche{}

\end{longtable}
% section registro_delle_modifiche (end)

\newpage
\thispagestyle{nopage}
\pagenumbering{roman}
\tableofcontents

\elencoFigure{}%
\elencoTabelle{}%

\newpage
\pagestyle{usual}
\pagenumbering{arabic}


\section{Informazioni generali}
\label{sec:info_generali}

\subsection{Informazioni incontro}
\label{sub:info_incontro}

\begin{itemize}
	\item \textbf{Luogo}: Applicazione \glossario{Zoom};
	\item \textbf{Data}: 2021-02-17;
	\item \textbf{Ora}: 11:30-12:30;
	\item \textbf{Partecipanti}:
	\begin{itemize}
		\item Contro Daniel Eduardo
		\item Fichera Jacopo
		\item Kostadinov Samuel
		\item Pagotto Manuel
		\item Paparazzo Giorgia
		\item Piccoli Gregorio
	\end{itemize}
\end{itemize}

\section{Ordine del giorno}%
\label{sec:ordine_del_giorno}

Di seguito vengono riportati i punti discussi durante l'incontro:
\begin{enumerate}
	\item \nameref{itm:1};
	\item \nameref{itm:2};
	\item \nameref{itm:3};
	\item \nameref{itm:4}.
\end{enumerate}


\subsection{Riferimenti}%
\label{sub:riferimenti}
\begin{itemize}
    \item \textsc{Analisi dei requisiti} v1.0.0;
	\item Esito I convocazione RR: \\
		\url{https://www.math.unipd.it/~tullio/IS-1/2020/Progetto/RR/SpaghettiCode.pdf}
\end{itemize}

\section{Resoconto}
\label{sec:resoconto}

Di seguito vengono riportate le decisioni concordate e gli esiti delle diverse discussioni con il proponente, oltre a 
delle note aggiuntive.

\subsection{Database supportati dal prodotto}
\label{itm:1}

\`{E} stato accordato che il prodotto supporter\`{a}, e permetter\`{a} di eseguire \glossario{query SQL}, sui \glossario{DBMS}: 
\glossario{MySQL}, \glossario{MariaDB},  \glossario{PostgreSQL}, \glossario{SQL Lite} e \glossario{Microsoft SQL Server}.

\subsubsection*{Documenti coinvolti}
Le modifiche relative alle decisioni prese saranno individuabili nel documento \textsc{Analisi dei Requisiti} a partire 
dalla versione v1.1.0.

\subsubsection*{Note Aggiuntive}
Sono state consigliate diverse tecnologie per interfacciarsi con i vari DBMS, tra cui \glossario{Apache drill}. Il gruppo
si prende la libert\`{a} di utilizzarle o meno, oltre che di aggiungere supporto ad ulteriori DBMS in futuro.

\subsection{Struttura dei file csv}
\label{itm:2}

Riguardo i file csv, che verranno forniti al sistema dall'utente, si possono fare le seguenti assunzioni: hanno la prima 
riga contenente i nomi delle varie dimensioni; i vari campi sono separati mediante il carattere ","; i campi stringa non 
contengono frasi bens\`{i} dati categorici di poche parole.

\subsection{Interazione con i database esterni}
\label{itm:3}
\`{E} stato deciso che l'interazione dell'utente con i database esterni si limiter\`{a} alla scelta di una delle query 
preconfigurate disponibili.

\subsubsection*{Documenti coinvolti}
La soluzione trovata sar\`{a} individuabile nel documento \textsc{Analisi dei Requisiti} a partire dalla versione v1.1.0.

\subsection{Use Case da implementare nel PoC}
\label{itm:4}
\`{E} stato concordato che nel PoC verranno implementati i casi d'uso:
\begin{itemize}
	\item UC1.2;
	\item UC2 ("via" UC2.2);
	\item UC3;
	\item UC4.2.1;
	\item UC4.2.2;
	\item UC4.2.4;
	\item UC4.2.5;
\end{itemize}
Il gruppo si tiene comunque la libert\`{a} di aggiungere l'implementazione di ulteriori casi d'uso al PoC qualora lo ritenesse valido.
\subsection{Note aggiuntive}
\label{sub:note}
Il proponente consiglia di effettuare i diversi calcoli di creazione e modifica dei grafici lato client e di prevedere un meccanismo di feedback,
per evitare disorientamento da parte dell'utente nel caso essi abbiano una durata consistente.
\end{document}
